\section{The Dempster--Shafer Reading}
\label{sec:ds}

Sections~\ref{sec:valuation}--\ref{sec:representation} independently
arrive at the same structure as Dempster--Shafer (DS)
theory~\cite{dempster1967,shafer1976}, including sub-additive complements, a
belief/plausibility interval, and mass functions.  
This section makes the identification exact in both statics and
dynamics.  In each case, it locates the precise point where excluded
middle makes a difference.

\subsection{Total Monotonicity via Constructive
Inclusion--Exclusion}

The $\neg\neg$-stable (\emph{regular}) elements of a frame form a
Boolean algebra with meet $\sqcap$ and join
$(a \sqcup b)^{\mathsf{cc}}$ (Glivenko).  Let us call the resulting 
algebra the
\emph{Booleanization} of this set.  
A DS \emph{belief function} is a capacity on a
Boolean algebra that is $n$-monotone for every $n$:
$\Bel(\bigvee_i A_i) \ge \sum_{\emptyset \neq T}
(-1)^{|T|+1} \Bel(\bigwedge_{i \in T} A_i)$.  We show $v$, restricted
to its Booleanization, is such a capacity for every $n$.  The
$2$-monotone case is immediate from modularity
(\texttt{two\_monotone}).  The full tower follows from the following 
result:

\begin{theorem}[inclusion--exclusion,
  \texttt{inclusion\_exclusion}]
\label{thm:incl-excl}
For every valuation on every frame $\Omega$, every finite family
$a : \iota \to \Omega$, and every finite $s \subseteq \iota$,
inclusion--exclusion holds with
equality for the frame join:
\[
  v\Bigl(\,\bigsqcup_{i \in s} a_i\Bigr)
  \;=\; \sum_{\emptyset \neq T \subseteq s} (-1)^{|T|+1}\,
        v\Bigl(\,\bigsqcap_{i \in T} a_i\Bigr).
\]
\end{theorem}

We state this in Lean additively, with the odd-cardinality half of the
signed sum set against the even-cardinality half, so that no
$\ENNReal$ subtraction occurs.  The proof is by induction with
modularity and frame distributivity.

Replacing the frame join $\bigsqcup_i a_i$ on the equality's other
side with its double negation $(\bigsqcup_i a_i)^{\mathsf{cc}}$, the
Booleanization's own join, yields the following inequality:

\begin{theorem}[total monotonicity, \texttt{infty\_monotone}]
\label{thm:infty}
For every valuation on every frame $\Omega$, every finite family
$a : \iota \to \Omega$, and every finite $s \subseteq \iota$:
\[
  \sum_{\emptyset \neq T \subseteq s} (-1)^{|T|+1}
  v\Bigl(\bigsqcap_{i \in T} a_i\Bigr)
  \;\le\;
  v\Bigl(\bigl(\bigsqcup_{i \in s} a_i\bigr)^{\mathsf{cc}}\Bigr).
\]
\end{theorem}

The Booleanization's join is given by
$(\cdot)^{\mathsf{cc}} \circ \sqcup$, and its meets of regular
elements are again regular. Interpreted in this way, 
Theorem~\ref{thm:infty} states exactly
the $\infty$-monotone tower: $v$ restricted to its Booleanization is
a belief function.

The proof of Theorem~\ref{thm:infty} from
Theorem~\ref{thm:incl-excl} follows from the monotonicity along
$a \le a^{\mathsf{cc}}$.  
Here, inclusion--exclusion does \emph{not} fail
constructively.  It is exact for the intuitionistic disjunction.
Applying the double-negation nucleus (a closure operator preserving
finite meets) $(\cdot)^{\mathsf{cc}}$ to the
join on the equality's right-hand side can only enlarge it. This 
follows from the fact that
$a \le a^{\mathsf{cc}}$ and $v$ are monotone, which turns the equality
into the inequality of Theorem~\ref{thm:infty}.  The amount by which
the right-hand side is enlarged is exactly the double-negation gap of
Section~\ref{sec:valuation}.

Additive probability does not become a
belief function by weakening an axiom, since modularity for $v$ is
never relaxed.  What changes is only the choice of join used to state the
tower: the frame join $\bigsqcup$ is replaced by its
double-negation closure $(\cdot)^{\mathsf{cc}} \circ \sqcup$.  A
belief function is thus the same additive valuation $v$, evaluated
against that closure rather than against the frame join.  With
$\Pl(a) = 1 - v(a^{\mathsf{c}})$, the belief--plausibility interval
has width exactly the slack
(\texttt{plausibility\_sub\_self}).  The DS ``ignorance'' is
boundary measure, not psychology.

\subsection{The Conditioning Characterization}

DS theory carries two update rules that coincide classically.
\emph{Geometric} conditioning renormalizes belief, and under the DS
dictionary it \emph{is} the localic posterior
$v(a \sqcap b)/v(b)$ of Section~\ref{sec:valuation}.

\begin{definition}[Dempster's rule of conditioning]
\label{def:dempster}
For a valuation $v$ on a frame $\Omega$ and $a, b \in \Omega$,
\emph{Dempster's rule of conditioning}~\cite{shafer1976} 
states that
\[
  v(a \,\|\, b) \;=\;
  \frac{v(a \sqcup b^{\mathsf{c}}) - v(b^{\mathsf{c}})}
       {1 - v(b^{\mathsf{c}})}.
\]
\end{definition}

\noindent
That is, the DS conditioning rule discards
the mass incompatible with the evidence and renormalizes by its
plausibility.
Both the geometric posterior of Section~\ref{sec:valuation} and
Dempster's rule of Definition~\ref{def:dempster} are again
valuations (\texttt{dempsterCond}).  Their modularity
is the frame distributivity
$(a \sqcup b^{\mathsf{c}}) \sqcap (a' \sqcup b^{\mathsf{c}})
= (a \sqcap a') \sqcup b^{\mathsf{c}}$.  Their disagreement is
given by the \emph{excluded-middle
gap} $\mathrm{emGap}(a,b) = v(a) - v(a \sqcap (b \sqcup
b^{\mathsf{c}}))$, i.e. the mass of $a$ outside $b$'s instance
of excluded middle:

\begin{lemma}[\texttt{sup\_compl\_decomp}]
For every valuation $v$ on every frame $\Omega$ and every
$a, b \in \Omega$,
$v(a \sqcup b^{\mathsf{c}}) = v(b^{\mathsf{c}}) + v(a \sqcap b) +
\mathrm{emGap}(a,b)$.
\end{lemma}

Dempster's numerator therefore exceeds Bayes's by exactly
$\mathrm{emGap}(a,b)$.  This gap depends on $a$, but whether it
vanishes at \emph{every} $a$ depends only on
the slack at the evidence $b$:

\begin{theorem}[the conditioning characterization,
  \texttt{dempsterCond\_eq\_\allowbreak condVal\_iff\_slack}]
\label{thm:cond-hinge}
For every valuation $v$ on every frame $\Omega$ and every
$b \in \Omega$ with $v(b) \neq 0$,
$\bigl(v(a \,\|\, b) = v(a \sqcap b)/v(b)$ for every
$a \in \Omega\bigr)$ iff $\slack(b) = 0$.
\end{theorem}

The hypothesis $\slack(b) = 0$ eliminates the gap
$\mathrm{emGap}(a,b)$ for every $a$ and makes the two normalizers
coincide, since $1 - v(b^{\mathsf{c}}) = v(b)$.
Conversely, testing at $a := b$, geometric conditioning gives
certainty $1$, while Dempster gives $v(b)/(1 - v(b^{\mathsf{c}}))$,
(which is $1$ only if the slack vanishes).
Theorem~\ref{thm:cond-hinge} is the valuation-relative analogue of
Theorem~\ref{thm:hinge}: both are iff statements about the same
proposition, $b \sqcup b^{\mathsf{c}} = \top$.  

We call the frame-level fact \emph{static} and the valuation-level
fact \emph{dynamic} because the former is a fixed structural property of
$\Omega$ with no notion of updating, while the latter characterizes
when a \emph{process} (conditioning $v$ on evidence $b$) behaves
classically.  The static characterization is an equation in the
frame itself, between $b \sqcup b^{\mathsf{c}}$ and $\top$, with no
valuation involved.  The dynamic characterization tests the same
proposition through $v$:
$\slack(b) = 0$ is exactly $v(b \sqcup b^{\mathsf{c}}) = v(\top)$.
Dempster and Bayes therefore coincide at $b$ whenever $v$ cannot
distinguish $b \sqcup b^{\mathsf{c}}$ from $\top$, even if the two
differ in the frame itself.  
In the measure model, Dempster and Bayes disagree precisely when the
evidence's boundary carries mass.
